\documentclass[11pt,a4paper]{article}

\usepackage[T1]{fontenc}
\usepackage[utf8]{inputenc}
\usepackage[french]{babel}
\usepackage{geometry}
\usepackage{hyperref}
\usepackage{longtable}
\usepackage{booktabs}
\usepackage{array}
\usepackage{enumitem}
\usepackage{xcolor}

\geometry{margin=2.5cm}
\hypersetup{
    colorlinks=true,
    linkcolor=blue,
    urlcolor=blue,
    pdftitle={SGEP --- Guide testeur API},
}

\title{SGEP --- Guide du testeur API\\
\large Parcours d'exploration Swagger / OpenAPI}
\author{Documentation SGEP}
\date{28 juin 2026}

\newcommand{\endpoint}[2]{\texttt{#1} & #2}

\begin{document}

\maketitle
\tableofcontents
\newpage

%==============================================================================
\section{Objectif}
%==============================================================================

Ce document guide une personne chargée de \textbf{tester et explorer l'API REST SGEP} via l'interface Swagger UI. Il décrit :
\begin{itemize}[noitemsep]
    \item comment accéder à la documentation interactive ;
    \item le flux d'authentification ;
    \item l'ordre recommandé pour couvrir tous les endpoints ;
    \item les champs auto-générés vs fournis par l'utilisateur ;
    \item les prérequis (rôles, données de test).
\end{itemize}

L'API est documentée avec \textbf{drf-spectacular} (OpenAPI~3). Les annotations source se trouvent dans \texttt{backend/core/openapi.py}.

%==============================================================================
\section{Accès à Swagger UI}
%==============================================================================

\subsection{URLs}

\begin{table}[h]
\centering
\begin{tabular}{ll}
\toprule
\textbf{URL} & \textbf{Description} \\
\midrule
\texttt{/api/docs/} & Interface Swagger UI (tests interactifs) \\
\texttt{/api/redoc/} & Documentation ReDoc (lecture) \\
\texttt{/api/schema/} & Schéma OpenAPI brut (JSON/YAML) \\
\texttt{/api/v1/} & Préfixe de toutes les routes métier \\
\bottomrule
\end{tabular}
\end{table}

En local : \texttt{http://localhost:8000/api/docs/}.\\
En déploiement : remplacer par l'URL du backend déployé (ex.\ Render).

\subsection{Générer le schéma en ligne de commande}

\begin{verbatim}
cd backend
python manage.py spectacular --color --file schema.yml
\end{verbatim}

%==============================================================================
\section{Authentification (première étape obligatoire)}
%==============================================================================

\subsection{Connexion}

\begin{enumerate}[noitemsep]
    \item Ouvrir le tag \textbf{Auth}.
    \item Exécuter \texttt{POST /api/v1/auth/login/} avec :
    \begin{verbatim}
    {
      "email": "admin@example.com",
      "password": "votre_mot_de_passe"
    }
    \end{verbatim}
    \item Copier \texttt{access\_token} dans la réponse.
    \item Cliquer sur \textbf{Authorize} (cadenas) en haut de Swagger UI.
    \item Saisir : \texttt{Bearer <access\_token>} (avec un espace après Bearer).
\end{enumerate}

\subsection{Rafraîchissement et déconnexion}

\begin{itemize}[noitemsep]
    \item \texttt{POST /api/v1/auth/refresh/} --- body \texttt{\{"refresh": "..."\}} --- renouvelle l'access token.
    \item \texttt{POST /api/v1/auth/logout/} --- blackliste le refresh token (auth requise).
    \item \texttt{POST /api/v1/auth/change-password/} --- ancien + nouveau mot de passe + confirmation.
\end{itemize}

\textbf{Note :} le login est limité à 5 tentatives / 10 minutes par IP.

%==============================================================================
\section{Rôles et permissions}
%==============================================================================

\begin{table}[h]
\centering
\begin{tabular}{lll}
\toprule
\textbf{Rôle} & \textbf{Endpoints principaux} & \textbf{Remarque} \\
\midrule
superadmin & Core, Classes, Students, Attendance, Grades, Report Cards, Admin & Gestion scolaire complète \\
comptable & Finance, Reports (exports) & Facturation et paiements \\
parent (actif) & Parent Portal (notes, absences, messages…) & Compte lié à un élève \\
parent (suspendu) & Profil, factures, date paiement prévue & Pas d'accès notes/absences \\
\bottomrule
\end{tabular}
\end{table}

Pour tester le portail parent, utiliser un compte parent créé à l'inscription d'un élève (champ \texttt{guardians}).

%==============================================================================
\section{Champs auto-générés vs fournis}
%==============================================================================

\subsection{Création d'un élève (\texttt{POST /api/v1/students/})}

Swagger marque le \textbf{matricule} comme champ requis du schéma. Comportement réel :
\begin{itemize}[noitemsep]
    \item Si \texttt{matricule} est \textbf{omis ou vide}, le serveur génère un identifiant unique (format \texttt{STU-AA-XXXXXX}).
    \item Si fourni, il doit être unique.
\end{itemize}

\textbf{Champs requis à la création :}
\texttt{first\_name}, \texttt{last\_name}, \texttt{birth\_date}, \texttt{birth\_place}, \texttt{gender}, \texttt{class\_id}, \texttt{academic\_year\_id}.

\textbf{Champs optionnels :}
\texttt{id\_number}, \texttt{school\_id}, \texttt{is\_active}, \texttt{medical}, \texttt{guardians}.

\textbf{Auto-générés côté serveur :}
\texttt{id}, \texttt{matricule} (si absent), \texttt{created\_at}, \texttt{updated\_at}, compte parent si \texttt{guardians} fourni.

\subsection{Factures}

\begin{itemize}[noitemsep]
    \item À l'\textbf{inscription} (\texttt{POST .../enroll/}) : seule la facture \texttt{INSC} est créée automatiquement.
    \item Les tranches \texttt{TR1}, \texttt{TR2}, \texttt{TR3} sont générées plus tard par le comptable via \texttt{POST /api/v1/finance/invoices/generate/}.
\end{itemize}

\subsection{Exports asynchrones}

Les endpoints d'export (élèves, absences, notes, finance) retournent \texttt{202} avec un \texttt{job\_id}. Suivre ensuite :
\begin{enumerate}[noitemsep]
    \item \texttt{GET /api/v1/reports/\{job\_id\}/status/}
    \item \texttt{GET /api/v1/reports/\{job\_id\}/download/} lorsque \texttt{status = completed}
\end{enumerate}

%==============================================================================
\section{Parcours de test recommandé}
%==============================================================================

Suivre cet ordre pour respecter les dépendances métier SGEP.

\subsection{Phase 0 --- Authentification}
\begin{enumerate}[noitemsep]
    \item \texttt{POST /api/v1/auth/login/}
    \item \texttt{GET /api/v1/admin/dashboard/} (vérifier accès superadmin)
\end{enumerate}

\subsection{Phase 1 --- Paramétrage (Core)}
\begin{enumerate}[noitemsep]
    \item \texttt{GET/POST /api/v1/schools/}
    \item \texttt{GET/PATCH /api/v1/schools/\{id\}/}
    \item \texttt{GET/POST /api/v1/academic-years/} --- activer l'année (\texttt{is\_active: true}) pour ouvrir les inscriptions
    \item \texttt{GET/PATCH /api/v1/academic-years/\{id\}/}
\end{enumerate}

\subsection{Phase 2 --- Structure pédagogique (Classes)}
\begin{enumerate}[noitemsep]
    \item \texttt{GET/POST /api/v1/classes/}
    \item \texttt{GET/PATCH /api/v1/classes/\{id\}/}
    \item \texttt{GET/POST /api/v1/subjects/}
    \item \texttt{GET/PATCH /api/v1/subjects/\{id\}/}
\end{enumerate}

\subsection{Phase 3 --- Élèves (Students)}
\begin{enumerate}[noitemsep]
    \item \texttt{POST /api/v1/students/} --- créer avec tuteurs (\texttt{guardians})
    \item \texttt{GET /api/v1/students/} --- pagination : \texttt{items}, \texttt{total}, \texttt{page}, \texttt{page\_size}
    \item \texttt{GET /api/v1/students/\{id\}/}
    \item \texttt{PATCH /api/v1/students/\{id\}/}
    \item \texttt{GET /api/v1/students/\{id\}/history/}
    \item \texttt{POST /api/v1/students/\{id\}/enroll/} --- inscription + facture INSC
    \item \texttt{POST /api/v1/students/\{id\}/promote/}
    \item \texttt{GET /api/v1/students/export/pdf/} puis suivi job
    \item \texttt{GET /api/v1/students/export/excel/} puis suivi job
    \item \texttt{DELETE /api/v1/students/\{id\}/} --- soft delete
\end{enumerate}

\subsection{Phase 4 --- Présences (Attendance)}
\begin{enumerate}[noitemsep]
    \item \texttt{GET/POST /api/v1/attendance/}
    \item \texttt{PUT /api/v1/attendance/\{id\}/}
    \item \texttt{POST /api/v1/attendance/\{id\}/justify/}
    \item \texttt{GET /api/v1/attendance/stats/} --- \texttt{class\_id}, \texttt{date\_from}, \texttt{date\_to} requis
    \item \texttt{GET /api/v1/attendance/export/}
\end{enumerate}

\subsection{Phase 5 --- Notes et bulletins (Grades / Report Cards)}
\begin{enumerate}[noitemsep]
    \item \texttt{GET/POST /api/v1/grades/}
    \item \texttt{PUT /api/v1/grades/\{id\}/}
    \item \texttt{POST /api/v1/grades/bulk/}
    \item \texttt{GET /api/v1/grades/export/results/}
    \item \texttt{POST /api/v1/report-cards/generate/}
    \item \texttt{GET /api/v1/report-cards/\{job\_id\}/status/} --- \texttt{pk} = ID du job, pas du bulletin
    \item \texttt{GET /api/v1/report-cards/\{id\}/download/}
    \item \texttt{POST /api/v1/report-cards/\{id\}/publish/}
\end{enumerate}

\subsection{Phase 6 --- Finance (Comptable)}
Se connecter avec un compte \textbf{comptable} ou superadmin.

\begin{enumerate}[noitemsep]
    \item \texttt{GET /api/v1/finance/dashboard/}
    \item \texttt{GET /api/v1/finance/invoices/}
    \item \texttt{POST /api/v1/finance/invoices/generate/} --- tranches TR1/TR2/TR3
    \item \texttt{GET /api/v1/finance/overdue/}
    \item \texttt{POST /api/v1/finance/payments/}
    \item \texttt{GET /api/v1/finance/payments/\{id\}/receipt/}
    \item \texttt{GET /api/v1/finance/export/excel/}
\end{enumerate}

\subsection{Phase 7 --- Portail parent}
Se connecter avec le compte parent.

\begin{enumerate}[noitemsep]
    \item \texttt{GET /api/v1/parent/me/}
    \item \texttt{GET /api/v1/parent/me/student/}
    \item \texttt{GET /api/v1/parent/me/grades/?period\_id=...}
    \item \texttt{GET /api/v1/parent/me/attendance/}
    \item \texttt{GET /api/v1/parent/me/report-cards/}
    \item \texttt{GET /api/v1/parent/me/report-cards/\{id\}/download/}
    \item \texttt{GET /api/v1/parent/me/invoices/}
    \item \texttt{PATCH /api/v1/parent/me/invoices/\{id\}/planned-payment-date/}
    \item \texttt{GET /api/v1/parent/me/invoices/\{payment\_id\}/receipt/} --- \texttt{pk} = ID paiement
    \item \texttt{GET/POST /api/v1/parent/me/messages/}
\end{enumerate}

\subsection{Phase 8 --- Jobs d'export (Reports)}
\begin{enumerate}[noitemsep]
    \item \texttt{GET /api/v1/reports/\{job\_id\}/status/}
    \item \texttt{GET /api/v1/reports/\{job\_id\}/download/}
\end{enumerate}

%==============================================================================
\section{Inventaire complet des endpoints}
%==============================================================================

\begin{longtable}{@{}p{6.8cm}p{1.2cm}p{6.5cm}@{}}
\toprule
\textbf{Chemin} & \textbf{Méthodes} & \textbf{Tag Swagger} \\
\midrule
\endhead
\endpoint{/api/v1/auth/login/}{POST} & Auth \\
\endpoint{/api/v1/auth/refresh/}{POST} & Auth \\
\endpoint{/api/v1/auth/logout/}{POST} & Auth \\
\endpoint{/api/v1/auth/change-password/}{POST} & Auth \\
\endpoint{/api/v1/admin/dashboard/}{GET} & Admin \\
\endpoint{/api/v1/schools/}{GET, POST} & Core \\
\endpoint{/api/v1/schools/\{id\}/}{GET, PATCH} & Core \\
\endpoint{/api/v1/academic-years/}{GET, POST} & Core \\
\endpoint{/api/v1/academic-years/\{id\}/}{GET, PATCH} & Core \\
\endpoint{/api/v1/classes/}{GET, POST} & Classes \\
\endpoint{/api/v1/classes/\{id\}/}{GET, PATCH} & Classes \\
\endpoint{/api/v1/subjects/}{GET, POST} & Classes \\
\endpoint{/api/v1/subjects/\{id\}/}{GET, PATCH} & Classes \\
\endpoint{/api/v1/students/}{GET, POST} & Students \\
\endpoint{/api/v1/students/\{id\}/}{GET, PATCH, DELETE} & Students \\
\endpoint{/api/v1/students/\{id\}/history/}{GET} & Students \\
\endpoint{/api/v1/students/\{id\}/enroll/}{POST} & Students \\
\endpoint{/api/v1/students/\{id\}/promote/}{POST} & Students \\
\endpoint{/api/v1/students/export/pdf/}{GET} & Students \\
\endpoint{/api/v1/students/export/excel/}{GET} & Students \\
\endpoint{/api/v1/attendance/}{GET, POST} & Attendance \\
\endpoint{/api/v1/attendance/\{id\}/}{PUT} & Attendance \\
\endpoint{/api/v1/attendance/\{id\}/justify/}{POST} & Attendance \\
\endpoint{/api/v1/attendance/stats/}{GET} & Attendance \\
\endpoint{/api/v1/attendance/export/}{GET} & Attendance \\
\endpoint{/api/v1/grades/}{GET, POST} & Grades \\
\endpoint{/api/v1/grades/\{id\}/}{PUT} & Grades \\
\endpoint{/api/v1/grades/bulk/}{POST} & Grades \\
\endpoint{/api/v1/grades/export/results/}{GET} & Grades \\
\endpoint{/api/v1/report-cards/generate/}{POST} & Report Cards \\
\endpoint{/api/v1/report-cards/\{id\}/status/}{GET} & Report Cards \\
\endpoint{/api/v1/report-cards/\{id\}/download/}{GET} & Report Cards \\
\endpoint{/api/v1/report-cards/\{id\}/publish/}{POST} & Report Cards \\
\endpoint{/api/v1/finance/invoices/}{GET} & Finance \\
\endpoint{/api/v1/finance/invoices/generate/}{POST} & Finance \\
\endpoint{/api/v1/finance/payments/}{POST} & Finance \\
\endpoint{/api/v1/finance/payments/\{id\}/receipt/}{GET} & Finance \\
\endpoint{/api/v1/finance/overdue/}{GET} & Finance \\
\endpoint{/api/v1/finance/dashboard/}{GET} & Finance \\
\endpoint{/api/v1/finance/export/excel/}{GET} & Finance \\
\endpoint{/api/v1/parent/me/}{GET} & Parent Portal \\
\endpoint{/api/v1/parent/me/student/}{GET} & Parent Portal \\
\endpoint{/api/v1/parent/me/grades/}{GET} & Parent Portal \\
\endpoint{/api/v1/parent/me/attendance/}{GET} & Parent Portal \\
\endpoint{/api/v1/parent/me/report-cards/}{GET} & Parent Portal \\
\endpoint{/api/v1/parent/me/report-cards/\{id\}/download/}{GET} & Parent Portal \\
\endpoint{/api/v1/parent/me/invoices/}{GET} & Parent Portal \\
\endpoint{/api/v1/parent/me/invoices/\{id\}/planned-payment-date/}{PATCH} & Parent Portal \\
\endpoint{/api/v1/parent/me/invoices/\{id\}/receipt/}{GET} & Parent Portal \\
\endpoint{/api/v1/parent/me/messages/}{GET, POST} & Parent Portal \\
\endpoint{/api/v1/reports/\{id\}/status/}{GET} & Reports \\
\endpoint{/api/v1/reports/\{id\}/download/}{GET} & Reports \\
\bottomrule
\end{longtable}

%==============================================================================
\section{Format des erreurs}
%==============================================================================

Les erreurs DRF sont encapsulées par le gestionnaire custom :
\begin{verbatim}
{
  "success": false,
  "error": {
    "message": "Description lisible",
    "details": { ... }
  }
}
\end{verbatim}

Codes courants : \texttt{400} validation, \texttt{401} non authentifié, \texttt{403} rôle insuffisant, \texttt{404} ressource absente, \texttt{409} conflit (ex.\ matricule dupliqué).

%==============================================================================
\section{Prérequis de test}
%==============================================================================

\begin{itemize}[noitemsep]
    \item Backend Django démarré avec variables d'environnement Appwrite configurées.
    \item Au moins un utilisateur \textbf{superadmin} et un \textbf{comptable} en base Appwrite.
    \item Clé \texttt{MEDICAL\_ENCRYPTION\_KEY} pour déchiffrer les données médicales côté admin.
    \item Redis + Celery pour les exports et bulletins asynchrones.
    \item Année scolaire active avant inscription d'élèves.
\end{itemize}

%==============================================================================
\section{Checklist de couverture}
%==============================================================================

Cocher chaque tag Swagger une fois tous ses endpoints testés :

\begin{itemize}[noitemsep]
    \item[$\square$] Auth (4 endpoints)
    \item[$\square$] Admin (1)
    \item[$\square$] Core (4 ressources)
    \item[$\square$] Classes (4 ressources)
    \item[$\square$] Students (8 chemins)
    \item[$\square$] Attendance (5 chemins)
    \item[$\square$] Grades (4 chemins)
    \item[$\square$] Report Cards (4 chemins)
    \item[$\square$] Finance (7 chemins)
    \item[$\square$] Parent Portal (10 chemins)
    \item[$\square$] Reports (2 chemins)
\end{itemize}

\textbf{Total : 52 opérations HTTP} documentées sur \textbf{46 chemins} sous \texttt{/api/v1/}.

\end{document}
